\section*{Chapter \ref{ch:ml:neural-networks-for-noisy-tabular-data} --- Neural Networks for Noisy Tabular Data}

\begin{solution}{exo:ml:neural-networks-for-noisy-tabular-data:1}
$40\times32 + 32 + 32\times16 + 16 + 16 + 1 = 1\,857$. With a first layer of 64: $40\times64 + 64 + 64\times32 + 32 +
32\times16 + 16 + 17 = 5\,249$.
\end{solution}

\begin{solution}{exo:ml:neural-networks-for-noisy-tabular-data:2}
$n \ge 8\times0.003^2/0.001^2 = 72$ seeds per architecture.
\end{solution}

\begin{solution}{exo:ml:neural-networks-for-noisy-tabular-data:3}
$160\,000/512 = 312.5$, so 313 steps; the chapter's $200\times500 = 100\,000$ rows give 196 steps an \omterm{def:ml:neural-networks-for-noisy-tabular-data:backprop}{epoch}.
\end{solution}

\begin{solution}{exo:ml:neural-networks-for-noisy-tabular-data:4}
The IC depends only on the ranking of the forecasts; R-squared also on their scale. Different seeds stop at different
\omterm{def:ml:neural-networks-for-noisy-tabular-data:backprop}{epochs} with forecasts spread more or less widely around a similar ranking, and the scale error enters R-squared
quadratically (\cref{prop:ml:why:ceiling}).
\end{solution}

\begin{solution}{exo:ml:neural-networks-for-noisy-tabular-data:5}
\omterm{def:ml:neural-networks-for-noisy-tabular-data:dropout}{Batch normalisation} standardises each unit with the batch's own statistics, which are noisy estimates on 512 rows of
mostly noise, and it couples the rows of a batch; a shallow network on ranked inputs gains nothing from better
conditioning in exchange. In a deep network normalisation keeps the layers' scales stable and usually makes training
possible at all; its noise is then a price worth paying.
\end{solution}

\begin{solution}{exo:ml:neural-networks-for-noisy-tabular-data:6}
The best of twenty validation scores is selected noise (\cref{prop:ml:validation:max}): the chosen seed is not better
out of sample than the others. Deploy the twenty-seed ensemble, which is.
\end{solution}

\begin{solution}{exo:ml:neural-networks-for-noisy-tabular-data:7}
\texttt{ensemble(n)}: one seed 0.0602, two 0.0625, five 0.0659, ten 0.0662. Two seeds buy $(0.0625 - 0.0602)/(0.0662 -
0.0602) = 38\%$ of the gain, five seeds 95\%.
\end{solution}

\begin{solution}{exo:ml:neural-networks-for-noisy-tabular-data:8}
With $z_1 = \mathrm{relu}(A_1x + c_1)$ and the forecast linear in $z_1$, the chain rule gives $\partial\ell_i/\partial A_1 =
(\partial\ell_i/\partial z_{1,i}\odot\mathbf 1\{A_1x_i + c_1 > 0\})\,x_i^\top$, the ReLU's derivative being the
indicator. The backward pass is one matrix-vector product per layer, as the forward pass is, plus elementwise products:
a small constant times the forward cost, independent of the number of parameters (Book~4, chapter~28).
\end{solution}

\begin{solution}{pb:ml:neural-networks-for-noisy-tabular-data:1}
\begin{enumerate}
\item 1\,857 and 5\,249.
\item After two to four \omterm{def:ml:neural-networks-for-noisy-tabular-data:backprop}{epochs}: the validation score peaks almost at once on data this noisy, and later \omterm{def:ml:neural-networks-for-noisy-tabular-data:backprop}{epochs} fit noise.
\item \omterm{def:ml:neural-networks-for-noisy-tabular-data:dropout}{Dropout} 0.2 raises the three-seed mean IC from 0.059 to 0.062; \omterm{def:ml:neural-networks-for-noisy-tabular-data:adam}{weight decay} 0.01 changes nothing; \omterm{def:ml:neural-networks-for-noisy-tabular-data:dropout}{batch
normalisation} lowers it to 0.052, \omterm{def:ml:neural-networks-for-noisy-tabular-data:dropout}{layer normalisation} to 0.056.
\item The seed standard deviation (0.0024) is comparable to the differences, so three seeds give standard errors of the
same size as the effects.
\item IC 0.057 to 0.065, R-squared 0.08\% to 0.30\%, net Sharpe ratio 1.37 to 1.78.
\item 0.0601 and 0.0619; gap 0.0019; pooled standard deviation 0.0024.
\item 14 per architecture.
\item $0.00244\sqrt2 = 0.0035$, almost twice the effect.
\item 0.066 each (the small ensemble: R-squared 0.32\%, net Sharpe ratio 1.73).
\item Averaging removes the seed-specific part of each member's error (\cref{prop:ml:trees:bagging}).
\item It disappears: both ensembles score 0.066.
\item The spread of the members' forecasts for each row, a measure of the model's own uncertainty (chapter~11).
\item A single network equals ridge (IC 0.060); the ensemble equals the lasso's IC (0.066) and trails the default boosted
trees (0.075, net Sharpe ratio 2.02).
\item \emph{Named result}: the seed standard deviation of the test IC, 0.0024, exceeds the architecture gap, 0.0019;
resolving it at two standard errors needs 14 seeds per architecture, and a ten-seed ensemble (0.066) beats both
architectures' single runs.
\item The boosted trees; if a network, its seed ensemble, never its best seed.
\item Architecture, every \omterm{def:ml:why-financial-machine-learning-is-different:sets}{hyperparameter}, the seeds, the validation block, the best \omterm{def:ml:neural-networks-for-noisy-tabular-data:backprop}{epoch} of each seed, the state hashes,
and the spread of the results.
\item More data lower the variance of each fit, so seeds matter less and larger networks can pay for themselves.
\item That the same code, data and seed produce the same weights; not that the result is good, nor that another seed would
agree.
\item The initial weights, the \omterm{def:ml:neural-networks-for-noisy-tabular-data:dropout}{dropout} masks and the order of the \omterm{def:ml:neural-networks-for-noisy-tabular-data:backprop}{mini-batches} (and, on accelerators, non-deterministic
kernels).
\item One draw from a distribution of models, whose spread must be measured before its mean is compared.
\end{enumerate}
\end{solution}

\begin{solution}{iq:ml:neural-networks-for-noisy-tabular-data:1}
The gradient of the loss with respect to every weight, by the chain rule applied from the output backwards (reverse-mode
differentiation). Its cost is a small multiple of one forward pass, because each layer's derivative is reused by all the
layers before it.
\iqlookfor{reverse mode and the cost argument.}
\end{solution}

\begin{solution}{iq:ml:neural-networks-for-noisy-tabular-data:2}
\omterm{def:ml:neural-networks-for-noisy-tabular-data:adam}{Adam} with an $\ell_2$ penalty adds $\lambda\theta$ to the gradient, which is then rescaled by \omterm{def:ml:neural-networks-for-noisy-tabular-data:adam}{Adam}'s per-parameter step, so
weights with large gradient variance are barely penalised; AdamW shrinks the weights directly, outside the adaptive step,
which restores \omterm{def:ml:neural-networks-for-noisy-tabular-data:adam}{weight decay}'s meaning.
\iqlookfor{the interaction between the penalty and the adaptive scaling.}
\end{solution}

\begin{solution}{iq:ml:neural-networks-for-noisy-tabular-data:3}
Train several seeds of each (\cref{prop:ml:nets:seeds} says how many), compare means with the standard error of the
difference, and compare the seed ensembles; check the gain on another period.
\iqlookfor{seeds as noise and a proper two-sample comparison.}
\end{solution}

\begin{solution}{iq:ml:neural-networks-for-noisy-tabular-data:4}
Seed every generator (torch, NumPy, the batch sampler), use one thread and deterministic algorithms
(\texttt{torch.use\_deterministic\_algorithms}),
fix library versions, and hash the weights. Different library versions, hardware, thread counts or BLAS back ends can
still change results in the last bits.
\iqlookfor{every source of randomness and of floating-point non-associativity.}
\end{solution}

\begin{solution}{iq:ml:neural-networks-for-noisy-tabular-data:5}
Trees: strong defaults, cheap tuning, robust to scaling, \omterm{def:ml:trees-and-boosting:monotone}{monotonic constraints}, fast inference, usually at least as good
on ranked characteristics. Networks: smooth interactions, embeddings, joint training with other objectives, benefits
from very large data, and ensembles that estimate their own uncertainty. On a monthly panel, start with trees and add a
network ensemble if it improves the stack.
\iqlookfor{a balanced argument tied to data size and use.}
\end{solution}

\begin{solution}{iq:ml:neural-networks-for-noisy-tabular-data:6}
Each network's error has a part common to all seeds (what the data and architecture allow) and a seed-specific part;
averaging removes the second, by $\rho\sigma^2 + (1-\rho)\sigma^2/B$. It stops helping when the seed-specific part is
gone, after a handful of seeds here: the members' remaining error is correlated.
\iqlookfor{the variance decomposition and its floor.}
\end{solution}
